%% ma: L11-B09-0007
%% lop: 11
%% chuong: 3
%% bai: 9
%% dang: 11.9.5
%% loai: TN4
%% muc: TH
%% dap_an: "C"
%% nguon: "Ảnh giáo viên gửi 24/10/2026 (M06): Kỳ thi thử đợt IV, Thi thử TN THPT 2027 (Biên soạn: Aimbet), Phần 1 Câu 6"
%% kien_thuc: "Bài 9 (tứ phân vị của mẫu số liệu ghép nhóm)"
%% trang_thai: chinh-thuc
\begin{ex}
Khảo sát cân nặng (đơn vị: kg) của $40$ trẻ sơ sinh tại một bệnh viện, người ta thu được bảng mẫu số liệu ghép nhóm như sau:
\begin{center}
\begin{tabular}{|c|c|c|c|c|}
\hline
\textbf{Cân nặng (đơn vị: kg)} & $[2{,}5;3)$ & $[3;3{,}5)$ & $[3{,}5;4)$ & $[4;4{,}5)$\\\hline
\textbf{Số trẻ sơ sinh} & $6$ & $16$ & $11$ & $7$\\\hline
\end{tabular}
\end{center}
Tứ phân vị thứ nhất $(Q_1)$ của bảng mẫu số liệu ghép nhóm đó bằng
\choice{$3{,}725$}{$3{,}375$}{\True $3{,}125$}{$3{,}625$}
\loigiai{$\dfrac n4=10$; tần số tích lũy: $6;\,22;\dots$ nên nhóm chứa $Q_1$ là $[3;3{,}5)$.
$Q_1=3+\dfrac{10-6}{16}\cdot0{,}5=3{,}125$.}
\end{ex}
